\chapter{Context and Objectives}

\vspace*{0.3cm}

Artificial Intelligence has become an important technology in modern applications,
especially with the development of language models capable of understanding
instructions and generating natural language responses. However, large language
models require significant computation, memory, and energy, which makes their
deployment difficult on constrained devices. Edge AI addresses this limitation by
executing models directly on local devices such as embedded systems, sensors, or
small computing boards. This improves latency, autonomy, and privacy, but it also
requires models that are compact and efficient \cite{edgeai}.

Small Language Models are therefore useful because they provide language
capabilities with a lower computational cost than large models. In this project,
the selected model is \texttt{google/flan-t5-small}, a compact instruction-tuned
model from the T5 family \cite{flant5model}. T5 represents tasks in a
text-to-text format, which is well suited for instruction-following data
\cite{t5}. Fine-tuning is used to adapt this pre-trained model to the prepared
dataset instead of training a model from scratch.

The objective of this work is to study the optimization of a Small Language Model
for Edge AI applications. The project prepares an instruction-tuning dataset,
fine-tunes \texttt{google/flan-t5-small}, and compares several optimization
algorithms: SGD, SGD with momentum, Adam, AdamW with clipping and scheduling,
L-BFGS, and Newton-CG. The comparison is not based only on validation loss, but
also on training latency, inference latency, throughput, and GPU memory usage.
These criteria are essential because an Edge AI model must be accurate, fast, and
resource-efficient.

The work also prepares a baseline for future compression. Techniques such as
quantization, pruning, and knowledge distillation are not fully implemented in
this report, but the optimizer comparison helps identify which fine-tuned model
would be the best candidate for later compression experiments.
